% 08_gravitational_waves.tex
% Gravitational Waves: GW propagation + GW detection as impedance antenna

\chapter{Gravitational Waves}
\label{ch:gravitational_waves}


\begin{objectivebox}
\begin{itemize}
    \item Derive gravitational wave propagation as lossless, trace-free, transverse impedance modulations of the LC vacuum lattice.
    \item Verify that GW strain amplitudes remain $\sim 28$ orders of magnitude below the dielectric hardware snap limit.
    \item Model LIGO/LISA interferometers as classical macroscopic impedance antennas responding to propagating lattice strain.
\end{itemize}
\end{objectivebox}

\section{GW Propagation as Impedance Modulation}
\label{sec:gw_propagation}

Gravitational waves are transverse inductive shear waves in the LC
lattice --- the same medium governed by the same operators. In accordance with the \textbf{Symmetric Scaling} axiom required to preserve a uniform optical vacuum, gravity scales both the local permittivity and permeability symmetrically via the refractive metric $n(r) = (1 - r_s/r)^{-1}$:
\begin{align}
    \varepsilon_{eff}(r) &= \varepsilon_0 \cdot n(r) \\
    \mu_{eff}(r) &= \mu_0 \cdot n(r)
\end{align}
where $r_s = 2GM/c^2$. Because both components scale proportionally, the macroscopic gravitational impedance remains invariant everywhere:
\begin{resultbox}{Invariant Gravitational Impedance}
\begin{equation}
    Z(r) = \sqrt{\frac{\mu_{eff}(r)}{\varepsilon_{eff}(r)}} = \sqrt{\frac{\mu_0 \cdot n}{\varepsilon_0 \cdot n}} \equiv Z_0
    \label{eq:Z_grav}
\end{equation}
\end{resultbox}
Therefore, the EM-transverse reflection coefficient $\Gamma_{EM}$ across any gravitational gradient evaluates identically to zero ($\Gamma_{EM} = 0$). While the local speed of light slows down near mass ($c_l = c/n$), the EM impedance matching remains perfect. In the weak-field regions traversed by detected signals, gravitational waves propagate without reflection or scattering. At the saturation boundary ($r_{sat} = 7GM/c^2$) the shear channel reflects totally ($\Gamma_{shear} = -1$, where $G_{shear} \to 0$), distinct from the EM channel which remains matched ($\Gamma_{EM} = 0$).

\begin{figure}[H]
    \centering
    \includegraphics[width=0.95\textwidth]{gw_schwarzschild_impedance.pdf}
    \caption{Schwarzschild impedance profile for a 30$\,M_\odot$ black hole,
    computed from the physics engine module \texttt{gw\_propagation.py}.
    \textbf{Top left:} Under Symmetric Gravity, both $\varepsilon$ and $\mu$ scale
    as $n(r)$, preserving $Z = Z_0$ everywhere outside the horizon.
    At the EM event horizon $r_s = 2GM/c^2$ the refractive index diverges
    ($n \to \infty$, $c_l \to 0$). The lattice phase transition ($S \to 0$,
    with the shear/bulk moduli collapsing) occurs at the deeper saturation
    boundary $r_{sat} = 7GM/c^2 = 3.5\,r_s$.
    \textbf{Top right:} Impedance $Z(r)/Z_0 = 1$ (invariant) in the propagation region.
    \textbf{Bottom left:} EM-channel reflection coefficient $\Gamma_{EM} = 0$
    (light is index-captured / transparent under symmetric gravity). The
    transverse-shear (GW) channel is distinct: at the saturation boundary
    $r_{sat} = 7GM/c^2$ the shear impedance $Z_{shear} \to 0$, giving
    $\Gamma_{shear} = -1$ (total reflection).
    \textbf{Bottom right:} Refractive index $n > 1$ near the mass
    (gravitational lensing = dielectric refraction).}
    \label{fig:gw_schwarzschild}
\end{figure}

\subsection{Lossless Propagation}

LIGO gravitational waves have strain
$h \sim 10^{-21}$.  The equivalent voltage per lattice cell is
$V_{GW} = h \cdot c \cdot \ell_{\text{node}} \cdot 2\pi f
\sim 10^{-19}\;\text{V}$, which is $\sim 2 \times 10^{-25}$ times smaller than
$V_{\text{snap}} \approx 511\;\text{kV}$.  Far below the saturation
threshold, the lattice acts as a perfect lossless linear transmission
line.  Gravitational waves propagate to cosmological distances without
dissipation --- matching LIGO observations exactly, with zero free
parameters.

\section{Orbital Damping: The Transmission Line Phase Slip}
\label{sec:orbital_damping}

While gravitational waves propagate losslessly through the linear regions of the LC lattice, their continuous generation requires real power to be extracted from the driving source. This is classically termed "Gravitational Radiation damping." Under the AVE framework, this mechanism is explicitly modeled as Topo-Kinematic \textbf{Transmission Line Phase Damping}.

\subsection{Hulse-Taylor Binary Pulsar (PSR B1913+16)}
Modeling the Hulse-Taylor binary pulsar as a macro-electrical orbital LC tank (\texttt{simulate\_\allowbreak binary\_\allowbreak lc\_\allowbreak damping.py}) reproduces the observed period decay rate. The quadrupole radiation reaction is re-expressed as transmission-line phase-slip damping; the $(v/c)^5$ quadrupole scaling and the Peters--Mathews coefficients are carried over from the standard radiation formula, so the result is a consistency reproduction rather than an independent first-principles derivation:

\begin{enumerate}
    \item \textbf{Reactive Tank Power ($Q$):} The gravitational and kinetic energy of the two $\approx 1.4 M_\odot$ neutron stars maps to a reactive resonant tank. Because Newtonian gravity holds an orthogonal V/I phase angle ($\cos(90^\circ) = 0$), the orbital energy oscillates losslessly.
    \item \textbf{Quadrupole Phase Slip ($\delta$):} The asymmetric quadrupole acceleration excites the surrounding elastic LC metric. Because lattice tension propagates at $c_0$, the differential acceleration wave acquires an inductive phase-retardation slip ($\delta$) over the metric span.
    % Rule-12 wording fix (2026-07-20, Grant Q1 ruling --- RULED loose-prose; Part 1 "loose
    % works"). Original step-3 wording, preserved verbatim here per KEEP-BOTH:
    %   "and real power radiates radially outwards into the bulk lattice network"
    % Grant ruled 2026-07-20 (Grant-verbatim three-parter, [sic]: "part 1: loose works /
    % Part 2: bank / part 3: rule it"; ruling content = the orchestrator's walk, ratified in
    % chat) that "bulk lattice network" is LOOSE PROSE for "the surrounding medium", NOT the
    % A1 / longitudinal bulk channel: the emission is SHEAR-channel (consistent with step-2's
    % $c_0$ propagation just above and the summary's exclusively-transverse statement). No
    % bulk/shear double-count exists. Wording corrected to "through the surrounding lattice as
    % transverse shear". Receipt chain in the (now-closed) warningbox below;
    % research/2026-07-20_q1-pulsar-hardening.md (merged #753).
    %
    % Rule-12 SCOPE REVISIT (2026-07-20, Grant F4 ruling --- manuscript-reconciliation lane).
    % Grant ruled the #766-review F4 routed question --> option (a), verbatim [sic]:
    %   "oh a for sure lets revisit that scope"
    % (_orchestration/2026-07-10_rulings-docket.md:2573 ENTRY 2026-07-20-masterq-tag-vocab-defs;
    % _orchestration/2026-07-20_pending-rulings-and-frontier-queue.md:26). Under Reading-A
    % (now standing physics per the Q1 revert), the Part-1 "loose works" shear-ONLY rewording is
    % INCOMPLETE: it correctly names the OBSERVED decay channel (shear, Reading-independent,
    % KEEP) but silently drops the framework's Reading-A bulk (longitudinal P-wave) admixture
    % (forbidden at O(1), pulsar-excluded). The Part-1 shear-only wording is PRESERVED here
    % verbatim per Rule 12 / KEEP-BOTH:
    %   "...real power radiates radially outwards through the surrounding lattice as transverse shear"
    % Item-3 is now restored to ADMIXTURE-HONEST wording: observed decay = shear (KEEP,
    % Reading-independent) + Reading-A bulk admixture = the LIVE pulsar exclusion + the OPEN
    % constituent-cage fork (see the manuscript-reconciliation warningbox below). The Part-1
    % identification of the OBSERVED channel as shear is NOT reverted (it survives untouched);
    % what is restored is the admixture honesty the shear-only phrasing lost under Reading-A.
    \item \textbf{Real Power Dissipation ($P_{real}$):} An electrical reactive tank experiencing a forced phase slip must bleed real power. The phase defect un-aligns the orthogonality, and real power radiates outward into the surrounding lattice ($P_{real} \approx Q \cdot \delta$). The \emph{observed} orbital-decay power is carried by the \textbf{transverse shear} channel --- the channel step-2 pins at $c_0$, the one the Peters--Mathews $(v/c)^5$ formula computes and pulsar timing measures; this observed-channel identification is Reading-independent and stands untouched. Under \textbf{Reading-A} (the standing physics as of 2026-07-20 --- see the warningboxes below), the framework \emph{additionally} predicts a bulk (longitudinal P-wave) admixture radiating on top of the shear decay at $O(1)$ coupling ($F_{\mathrm{bulk}}/F_{\mathrm{shear}} \approx 0.03$--$0.12$); that admixture is \emph{not} part of the observed shear decay, and its standing status is the \textbf{LIVE} pulsar exclusion plus the \textbf{OPEN} constituent-cage fork flagged in the reconciliation box below.
    \item \textbf{Orbital Decay ($\dot{P}_b$):} Subjecting this real power drain over the empirical Hulse-Taylor orbital bounds algebraically evaluates an exact period decay rate of $\mathbf{-2.402 \times 10^{-12}\text{ s/s}}$.
\end{enumerate}

The circuit model reproduces the observed orbital decay to within $\sim 2\%$, showing that the general-relativistic quadrupole result can be cast as transmission-line phase-retardation across the LC lattice.

\begin{figure}[htbp]
    \centering
    \includegraphics[width=0.85\textwidth]{fig_hulse_taylor_phase_slip.png}
    \caption{\textbf{Quadrupole phase-slip damping of the binary orbital tank.} The blue branch is the radiated power \textbf{delivered to the far-field port}, $P_{\mathrm{rad}} = Q\,\delta$, against the reactive tank power $Q$ of the orbital LC tank; the grey dashed line is the reactive ceiling $P = Q$; the star marks PSR~B1913+16 at its empirical orbital bounds. \textbf{Transfer-cost framing, and it is load-bearing:} the ordinate is power \emph{delivered}, \textbf{not} power ``dissipated'' --- nothing in the binary is a resistor. Per the transfer-cost theorem (\kbleaf{ave-kb/common/transfer-cost-theorem.md}, \texttt{clm-xfrcst}), within-system reactive exchange is free however many channels it crosses; an arrow appears only where the transfer crosses the \emph{system boundary} through a counting port, and a real port at the boundary is a \textbf{delivery terminal} --- the substrate does not heat, the energy leaves intact into whatever counts it. \emph{Scope:} the $(v/c)^5$ quadrupole scaling and the Peters--Mathews coefficients are carried over from the standard radiation formula (see the derivation above), so this is a \textbf{consistency reproduction} cast in circuit variables, not an independent first-principles derivation of the decay rate.}
    \label{fig:hulse_taylor_phase_slip}
\end{figure}

% Rule-12 closure (2026-07-20, Grant Q1 ruling): the "unresolved channel attribution"
% warningbox is CLOSED with a dated RESOLVED-with-receipt lead-in; the ORIGINAL warningbox
% text is PRESERVED VERBATIM below the lead-in per KEEP-BOTH (do not delete it). Grant ruled
% 2026-07-20 (Grant-verbatim three-parter, [sic]: "part 1: loose works / Part 2: bank /
% part 3: rule it"; ruling content = the orchestrator's walk, ratified in chat): step-3's
% "bulk lattice network" is loose prose for the surrounding medium, NOT the A1 bulk channel;
% the emission is shear-channel; no bulk/shear double-count exists. Receipts: the pulsar
% exclusion (merged #753, research/2026-07-20_q1-pulsar-hardening.md) + step-2's own c_0.
\begin{warningbox}[RESOLVED 2026-07-20 (Grant Q1 ruling) --- channel attribution (formerly: unresolved channel attribution)]
\textbf{RESOLVED 2026-07-20 (Grant Q1 ruling --- Part~1 ``loose works''; ruling content = the
orchestrator's walk, ratified in chat).} The channel-attribution tension flagged below is
\emph{closed}. Step-3's ``radiates \ldots{} into the bulk lattice network'' was \textbf{loose
prose} for ``the surrounding medium'', \textbf{not} the A1 / longitudinal bulk channel: the
binary's orbital-decay power is emitted through the \textbf{transverse shear} channel --- the
one channel the summary box below and the KB canon already name --- so \textbf{no bulk/shear
double-count exists}. Step-3's wording is corrected to ``through the surrounding lattice as
transverse shear'' (Rule-12 note at item~3 above). \textbf{Receipts:} (i) the pulsar-timing
exclusion --- an \emph{independent} radiative bulk port with $O(1)$ coupling is excluded at
$9$--$110\sigma$ by Hulse-Taylor and by $100$--$1400\times$ the double-pulsar bound
(\texttt{research/2026-07-20\_q1-pulsar-hardening.md}, merged \#753), which forbids the
double-count directly; (ii) step-2 of this same derivation (item~2 above) already pins the
decay wave's propagation at $c_0$ --- the \emph{shear} speed --- internal textual evidence
that step-3 meant the medium, not the bulk channel. The Reading-A radiative-port exclusion is
banked in \texttt{ave-kb/common/genesis-chord-falsification-ledger.md}; the Reading-B
standing ruling (no far-field bulk radiative port for gravitating sources; the halo is the
complete reactive-near-field story) is \textbf{CONDITIONAL} on the B-mechanism lane deriving
the closure (if that lane lands none-derives, the Reading-A exclusion goes live).

\medskip
\noindent\emph{Original current-status flag preserved verbatim below per Rule~12 / KEEP-BOTH
(the wording ``left unedited pending it'' records the pre-ruling state; item~3 is now
corrected):}

\medskip
\noindent
The phrase ``real power radiates radially outwards into the bulk lattice
network'' in item~3 above (the $P_{real}$ step) attributes the radiated
power to the \emph{bulk / longitudinal} channel. This is in tension with (i)
the trace-free, transverse summary of this same chapter (the summary box
below: ``Gravitational waves propagate exclusively as lossless, trace-free,
transverse impedance modulations''); (ii) the KB canon, which assigns
gravitational waves to the \emph{transverse shear} channel
(\kbleaf{vol3/gravity/ch08-gravitational-waves/gw-propagation-lossless.md}
``transverse inductive shear waves'';
\kbleaf{vol9/ch4-dc-electrical-characteristics/three-channel-impedances.md},
the ``Shear / GW'' row). The same contradiction was independently surfaced by
the vacuum-impedance-probe Phase-A feasibility result (which flags that the
corpus consistently assigns GW to the shear channel while the bulk /
V-sector is the genuinely \emph{un}-instrumented channel) and by the
2026-07-03 longitudinal-energy-budget exposure audit. Channel attribution
for the binary-pulsar phase-slip damping \textbf{awaits Grant adjudication}
and is not resolved here; the wording of item~3 is left unedited pending it.
A candidate route to a decisive resolution is an engine channel-decomposition
of the binary's radiated power (shear vs.\ bulk), which would fork this from
a prose ambiguity to a computable quantity. See the synthesis at
\texttt{research/2026-07-03\_historical-exposure-audits\_synthesis.md}.
\end{warningbox}

% Q1 REVERT flag (2026-07-20). Additive status warningbox --- physics prose NOT
% rewritten (flag-don't-fix / Rule 12), the RESOLVED box above is NOT edited. The
% RESOLVED closure above rested in part on the RULED-CONDITIONAL Reading-B ruling
% (no far-field bulk radiative port), which that same box flagged as CONDITIONAL
% on the B-mechanism lane. That lane landed NONE-DERIVES at review grade (#761
% merged, main @ caa51c17; research/2026-07-20_mechanical-commonmode-derivation_result.md),
% so per port-register.md:87's own clause Reading-A goes live. The Part-1
% "loose works" fix at item 3 (step-3 = transverse shear) is a SEPARATE ruling
% and is NOT reverted; what reverts is the settled-ness of "the bulk channel does
% not radiate." Routed to the manuscript-reconciliation program.
\begin{warningbox}[Q1 REVERT 2026-07-20 --- bulk radiative port DERIVED-OPEN + pulsar-excluded (LIVE open contradiction, routed)]
\textbf{Status change 2026-07-20 (Q1 revert --- supersedes the ``no bulk/shear
double-count exists'' \emph{basis} of the RESOLVED box above; the RESOLVED box is
NOT edited; ruling content = the orchestrator's walk, ratified in chat).} The
RESOLVED closure above rested in part on the then-standing \emph{RULED-CONDITIONAL
Reading-B} ruling (``no far-field bulk radiative port for gravitating sources; the
halo is the complete reactive-near-field story'') --- explicitly conditional, in
that same box, on the B-mechanism lane deriving the closure. \textbf{That lane has
now landed NONE-DERIVES} at review grade (\#761 merged, main @ \texttt{caa51c17};
\texttt{research/2026-07-20\_mechanical-commonmode-derivation\_result.md}, \S5):
the A1/dilatation mass channel rides a genuine gapless P-branch, a gravitating
quadrupole drives it at quadrupole order, and the EM Gauss-kill is
\emph{structurally blocked} from the mechanical $\nabla\!\cdot\!u$ by the bulk
restoring force $K\neq0$ --- so AVE's own elastic-solid vacuum \emph{forces} a
radiating bulk (P-wave) channel (Reading~A). Per the port-register Q1 row's own
clause (\texttt{port-register.md:87}, verbatim [sic]: ``if NONE-DERIVES this
ruling REVERTS and the banked Reading-A exclusion (falsification ledger;
9--110$\sigma$ pulsar) goes live''), \textbf{Reading-B has REVERTED and the
Reading-A exclusion is now LIVE against the framework}: an independent far-field
bulk radiative port at $O(1)$ coupling is excluded at $9$--$110\sigma$
(Hulse-Taylor) and by $100$--$1400\times$ the double-pulsar bound.

\medskip
\noindent\textbf{Consequence for this chapter (LIVE open contradiction, routed
--- physics prose NOT rewritten here).} This chapter's Hulse-Taylor macro-LC
reproduction (\S\ref{sec:orbital_damping}, items~1--4) rides the standard
\emph{shear}-quadrupole (Peters--Mathews $(v/c)^5$) radiation formula and
attributes the \emph{entire} orbital decay to the transverse shear channel. With
Reading-A now standing, the framework additionally predicts a \emph{bulk}
radiative admixture ($F_{\mathrm{bulk}}/F_{\mathrm{shear}}\approx0.03$--$0.12$ at
$O(1)$ coupling) \emph{on top of} the shear decay --- the q1-hardening~\S4
bulk/shear double-count --- which binary-pulsar timing \emph{excludes}. So the
chapter's gravitational-radiation story now carries a \textbf{LIVE open
contradiction}: either the bulk channel radiates too (pulsar-excluded,
kill-class) or an as-yet \emph{underived} suppression is owed (the Reading-B
suppression the \#761 derivation found NONE-DERIVES). The Part-1 ``loose works''
correction at item~3 (step-3 $=$ transverse shear) is a \emph{separate} ruling and
is \textbf{not} reverted; what reverts is the settled-ness of ``the bulk channel
does not radiate.'' This contradiction is \textbf{routed to the
manuscript-reconciliation program}, not resolved here (flag-don't-fix / Rule~12).
\textbf{Receipts:} the \#761 result \S5 CONSEQUENCE
(\texttt{research/2026-07-20\_mechanical-commonmode-derivation\_result.md}); the
reverted port-register Q1 row
(\kbleaf{common/port-register.md} \S3); the falsification-ledger entry
\texttt{q1-reading-A-radiative-bulk-port} (promoted LIVE 2026-07-20,
\kbleaf{common/genesis-chord-falsification-ledger.md}).

\medskip
\noindent\emph{[Superseded-in-part 2026-07-20 by the manuscript-reconciliation box below.]} The
``physics prose NOT rewritten here / routed'' state recorded in this box is the state at
\emph{Q1-revert-execution} time; the manuscript-reconciliation lane has since \textbf{discharged
that routing} --- item-3 is now reworded admixture-honest (Grant's F4$\to$(a) ruling), and the
constituent-cage fork is stated as its honest \textbf{OPEN} state. This box is preserved verbatim
(Rule~12); read the reconciliation box below as current.
\end{warningbox}

% Manuscript-reconciliation lane (2026-07-20). Additive dated warningbox discharging the routing
% of the Q1 REVERT box above and reconciling the chapter to current canon (the honesty-lag
% discipline): (i) item-3 reworded admixture-honest per Grant's F4 --> (a) ruling; (ii) the
% #770 constituent-cage fork stated as its honest OPEN state (deep-rail frozen-first program
% RUNNING concurrently --- outcome NOT presumed). Both prior boxes preserved verbatim (Rule 12).
\begin{warningbox}[Manuscript-reconciliation 2026-07-20 --- item-3 admixture-honest (F4$\to$(a)) + the constituent-cage fork is OPEN (deep-rail program running)]
\textbf{Reconciliation to current canon (2026-07-20, manuscript-reconciliation lane; the
honesty-lag discipline --- the chapter must not lag the KB truth-source).} This box discharges the
routing of the Q1 REVERT box above; both prior boxes are preserved verbatim (Rule~12).

\medskip
\noindent\textbf{(1) Item-3 reworded admixture-honest (Grant F4 ruling $\to$ option (a)).} Grant
ruled the \#766-review F4 routed question $\to$ option~(a), verbatim [sic]: \emph{``oh a for sure
lets revisit that scope''} (\texttt{\_orchestration/2026-07-10\_rulings-docket.md:2573};
\texttt{\_orchestration/2026-07-20\_pending-rulings-and-frontier-queue.md:26}). Under Reading-A the
Part-1 ``loose works'' \emph{shear-only} rewording of item-3 was \emph{incomplete}: it correctly
names the \textbf{observed} decay channel as transverse shear (Reading-independent --- what the
Peters--Mathews $(v/c)^5$ formula computes and pulsar timing measures; this survives untouched),
but it silently dropped the framework's Reading-A \textbf{bulk (longitudinal P-wave) admixture}.
Item-3 above now carries both: observed decay $=$ shear (KEEP) $+$ the forbidden-at-$O(1)$ bulk
admixture ($F_{\mathrm{bulk}}/F_{\mathrm{shear}} \approx 0.03$--$0.12$). The Part-1 step-3$=$shear
ruling for the observed channel is \textbf{not} reverted.

\medskip
\noindent\textbf{(2) The bulk admixture's standing status is a LIVE exclusion \emph{plus} an OPEN
fork --- stated honestly, not resolved.} Two distinct things are true simultaneously and the
chapter states both:
\begin{itemize}
    \item \textbf{LIVE exclusion (settled negative).} Under Reading-A --- now the standing physics
    (Q1 revert; \#761 make-or-break mechanical $\nabla\!\cdot\!u$ derivation returned
    \textbf{NONE-DERIVES}, \texttt{research/2026-07-20\_mechanical-commonmode-derivation\_result.md}
    \S5) --- an \emph{independent far-field bulk radiative port at $O(1)$ coupling} is excluded by
    binary-pulsar timing at $9$--$110\sigma$ (Hulse-Taylor) and by $100$--$1400\times$ the
    double-pulsar bound (falsification-ledger entry
    \texttt{q1-reading-A-radiative-bulk-port}, promoted LIVE 2026-07-20). A Reading-A bulk
    admixture \emph{as an independent radiative port} is a closed-negative against the framework
    (Rule~11 honest closure).
    \item \textbf{OPEN constituent-cage fork (running --- outcome NOT presumed).} Whether the
    \emph{deeper} constituent structure evades that exclusion is \textbf{re-opened}, not resolved.
    The AVE mass texture is an ensemble of $\sim\!10^{57}$ constituent solitons, each canonically
    carrying a bulk-only $\Gamma_{\mathrm{bulk}}=-1$ knot-core cage. The question: does that
    ensemble present to the $2\Omega$ far field the coarse-grained \emph{uncaged} mass texture
    (which \textbf{radiates} at the q1 partition $\Rightarrow$ confirms the kill one level deeper),
    or do the per-core $\Gamma_{\mathrm{bulk}}=-1$ boundaries \textbf{image-cancel} the long-wave
    compression moment (Lloyd-mirror class $\Rightarrow$ a derived suppression that could later
    re-open Reading-B)? The maximum-stakes review of \#770
    (\texttt{research/2026-07-20\_constituent-cage-ensemble\_result.md}, \S0/\S7) \textbf{refuted}
    the earlier ``ensemble radiates, fork closes'' bank: at deep rail the canon channel-asymmetric
    bulk-only wall \emph{is} constitutively realizable ($\Gamma_{\mathrm{bulk}}\!\to\!-1$ with the
    shear channel finite), so the image-cancel (suppression) route is \textbf{live again}. The
    verdict is re-banked \textbf{WALL-CLASS / RAIL-DEPTH-CONDITIONAL $\to$ REOPENED}; the decisive
    adjudicator is the \textbf{deep-rail frozen-first program} (\#770~\S8.2), which is
    \textbf{a registered owed follow-on (\#770 \S 8.2), dispatched 2026-07-20}. \emph{Its outcome is not presumed here} --- the chapter states
    the fork as open, in either direction.
\end{itemize}

\medskip
\noindent\textbf{Net standing state for this chapter.} The Hulse-Taylor macro-LC reproduction
(\S\ref{sec:orbital_damping}) is a consistency reproduction of the \emph{observed shear-quadrupole}
decay and stands. The framework's Reading-A \emph{bulk admixture} is a live pulsar-excluded
prediction \emph{at the naive independent-port reading}, with the constituent-cage suppression fork
genuinely OPEN pending the deep-rail program --- so the chapter neither claims a clean gravitational
chord here nor pretends the bulk sector is settled. \textbf{Receipts:} \#761 result \S5
(\texttt{research/2026-07-20\_mechanical-commonmode-derivation\_result.md}); the \#770 REOPENED
verdict + \S8.2 deep-rail program
(\texttt{research/2026-07-20\_constituent-cage-ensemble\_result.md}); the reverted port-register
Q1 row (\kbleaf{common/port-register.md}~\S3); the falsification-ledger entry
\texttt{q1-reading-A-radiative-bulk-port} (\kbleaf{common/genesis-chord-falsification-ledger.md}).
\end{warningbox}

\section{GW Detection: The Impedance Antenna}
\label{sec:gw_detection}

A gravitational wave detector is an impedance antenna.  LIGO's 4~km
Fabry--P\'erot arms form resonant cavities in the LC vacuum, where
each light bounce amplifies the GW-induced impedance modulation.
The passing GW strain $h$ perturbs the local vacuum impedance:
\begin{resultbox}{GW-Induced Shear-Impedance Perturbation}
\begin{equation}
    \delta Z_{shear} = Z_{shear} \cdot h
\end{equation}
\end{resultbox}
where $Z_{shear} = \rho\,c_{shear}$ is the shear (transverse-mechanical) impedance. The EM characteristic impedance $Z_0 \equiv Z_{EM}$ is the matched radiative port and is invariant under symmetric gravity; it is a distinct channel.
The accumulated phase shift after $N$ bounces is:
\begin{resultbox}{Fabry-P\'erot Accumulated Phase Shift}
\begin{equation}
    \Delta\phi = \frac{2\pi f_{GW}}{c} \cdot L \cdot N \cdot h
\end{equation}
\end{resultbox}
For LIGO ($L=4$~km, $N=280$, $h=10^{-21}$, $f=100$~Hz),
$\Delta\phi \approx 2.3 \times 10^{-21}$~rad --- resolved via
homodyne readout against 750~kW circulating laser power.

The strain sensitivity is bounded by two quantum noise sources:
shot noise (phase) and radiation pressure (amplitude), whose
geometric mean is the Standard Quantum Limit:
\begin{resultbox}{Standard Quantum Limit (Strain)}
\begin{equation}
    h_{SQL}(f) = \sqrt{h_{shot}^2 + h_{RP}^2}
\end{equation}
\end{resultbox}

\begin{figure}[H]
    \centering
    \includegraphics[width=0.95\textwidth]{gw_detector_sensitivity.pdf}
    \caption{\textbf{Left:} Strain sensitivity curves for LIGO (4~km),
    LISA (2.5~Gm), and Einstein Telescope (10~km), modeled as
    impedance antennas. GW150914 ($h\sim10^{-21}$) is marked.
    \textbf{Right:} LIGO noise budget decomposition: shot noise
    dominates at high frequency, radiation pressure at low frequency,
    with the SQL minimum determining peak sensitivity.}
    \label{fig:gw_detector}
\end{figure}

\noindent
The lattice voltage ratio for LIGO GW is:
\begin{resultbox}{LIGO GW Saturation Ratio}
\begin{equation}
    \frac{V_{GW}}{V_{\text{snap}}} \approx 2 \times 10^{-25}
\end{equation}
\end{resultbox}
Twenty-five orders of magnitude below saturation.  The vacuum is
a \textit{perfect} lossless transmission line for gravitational waves,
exactly as observed.

\begin{summarybox}
\begin{itemize}
    \item Gravitational waves propagate exclusively as lossless, trace-free, transverse impedance modulations of the macroscopic LC vacuum lattice.
    \item The spatial vacuum operates as a perfect transmission line because GW voltage amplitudes remain $\sim 28$ orders of magnitude below the dielectric hardware snap limit.
    \item Interferometric GW detectors (LIGO, LISA) are modeled natively as classical macroscopic impedance antennas responding perfectly to propagating lattice strain ($\Delta\phi$).
\end{itemize}
\end{summarybox}

\begin{exercisebox}
\begin{enumerate}
    \item Show mathematically that the $V_{GW}/V_{snap}$ ratio for a standard LIGO detection guarantees perfectly linear wave propagation free from dielectric non-linearities.
    \item Calculate the differential node displacement (strain length) along a 4 km LC network segment perturbed by a target $h \sim 10^{-21}$ gravitational wave amplitude.
\end{enumerate}
\end{exercisebox}
